% Part: computability
% Chapter: computability-theory
% Section: ce-sets

\documentclass[../../../include/open-logic-section]{subfiles}

\begin{document}

\olfileid[fr]{cmp}{thy}{ces}
\olsection{Ensembles effectivement énumérables}

\begin{defn}
Un ensemble est \emph{effectivement énumérable}
s'il est vide ou s'il est l'image d'une fonction
calculable totale, définie sur $\Nat$.
\end{defn}

\begin{history}
On dit aussi \emph{récursivement énumérable}.
C'est la terminologie d'origine; les deux
appellations restent courantes aujourd'hui,
de même que les abréviations anglaises
« c.e. » et « r.e. ».
\end{history}

\begin{explain}
Il convient de réfléchir au sens de cette
définition et à la pertinence du terme.
Si $S$ est l'image d'une fonction
calculable~$f$, alors
\[
S = \{ f(0), f(1), f(2), \dots \},
\]
de sorte que $f$ « énumère » les éléments
de~$S$. La définition n'exige pas que
$f$ soit croissante : les éléments n'ont
pas besoin d'être énumérés dans l'ordre.
Elle n'exige pas non plus que $f$ soit
injective : les répétitions dans la
suite $f(0),f(1),f(2),\dots{}$ sont permises.
Par exemple, la fonction constante
$f(x)=0$ énumère l'ensemble $\{0\}$.
\end{explain}

Tout ensemble calculable est effectivement
énumérable. En effet, supposons $S$
calculable. Si $S$ est vide, il satisfait
la définition. Sinon, choisissons un
élément~$a$ de~$S$ et définissons
$f$ par
\[
f(x) =
\begin{cases}
x & \text{si $\Char{S}(x)=1$}\\
a & \text{sinon.}
\end{cases}
\]
La fonction $f$ est calculable et son
image est exactement~$S$.

\end{document}
